\chapter{Modelarea bazei de date și arhitectura sistemului}

\section{Introducere}

Aici arăt cum arată baza de date din spatele aplicației: ce entități am definit,
cum se leagă între ele și cum am împărțit codul pe module.

\section{Modelarea bazei de date}

Baza de date este construită pe un model relațional~\cite{elmasri-navathe},
gestionat prin Entity Framework Core. Pe lângă entitățile proprii aplicației, schema
include și tabelele standard generate de \emph{ASP.NET Identity}
(\texttt{AspNetUsers}, \texttt{AspNetRoles} etc.), entitatea \texttt{ApplicationUser}
extinzând clasa \texttt{IdentityUser}.

\subsection{Entități principale}

Denumirile claselor păstrează terminologia engleză folosită la implementare
(\texttt{Topic} pentru materie, \texttt{HireRequest} pentru cererea de colaborare),
în timp ce textul lucrării și interfața folosesc termenii românești.

\begin{description}
    \item[ApplicationUser] reprezintă un cont de utilizator; extinde \texttt{IdentityUser}
    cu câmpurile \texttt{Bio}, \texttt{Role} (\texttt{"Teach"} sau \texttt{"Learn"}),
    \texttt{EducationLevel}, \texttt{TeachingLevel}, \texttt{Age}, \texttt{Gender},
    \texttt{SelectedTopics} (identificatorii materiilor alese, separați prin virgulă),
    \texttt{AvatarPath} și \texttt{BannedUntil} (data până la care contul e blocat).
    \item[Message] reprezintă un mesaj din conversația dintre doi utilizatori: conține
    \texttt{SenderId}, \texttt{ReceiverId}, \texttt{Content}, \texttt{Timestamp},
    \texttt{IsRead}, precum și câmpurile opționale \texttt{FilePath}, \texttt{FileName}
    și \texttt{FileSize}, populate atunci când mesajul conține un fișier atașat.
    \item[HireRequest] reprezintă o cerere de colaborare trimisă de un elev unui profesor:
    conține \texttt{RequesterId}, \texttt{TargetId}, \texttt{IsAccepted} și
    \texttt{RequestedAt}.
    \item[Topic] reprezintă o materie (de exemplu, Matematică, Programare);
    conține \texttt{Name} și \texttt{Description}.
    \item[OtpCode] reprezintă un cod de unică folosință generat pentru verificarea
    email-ului sau confirmarea unei modificări sensibile a profilului; conține
    \texttt{UserId}, \texttt{Code} și \texttt{ExpirationTime}.
\end{description}

\subsection{Relații importante}

Figura~\ref{fig:er-diagram} prezintă relațiile principale dintre entități.

\begin{figure}[h]
\centering
\begin{tikzpicture}[
    box/.style={draw, rectangle, rounded corners=2pt, minimum width=2.6cm, minimum height=0.8cm, align=center, font=\small},
    >=Stealth
]
\node[box] (user) {ApplicationUser};
\node[box, right=3.6cm of user, yshift=2.4cm] (message) {Message};
\node[box, right=3.6cm of user, yshift=0.8cm] (hire) {HireRequest};
\node[box, right=3.6cm of user, yshift=-0.8cm] (otp) {OtpCode};
\node[box, right=3.6cm of user, yshift=-2.4cm] (topic) {Topic};
\node[box, above=1.8cm of user] (role) {AspNetRoles};

\draw[->] (user) -- (message) node[midway, above, sloped, font=\tiny] {1..n (Sender / Receiver)};
\draw[->] (user) -- (hire) node[midway, above, sloped, font=\tiny] {1..n (Requester / Target)};
\draw[->] (user) -- (otp) node[midway, above, sloped, font=\tiny] {1..n};
\draw[->] (user) -- (topic) node[midway, above, sloped, font=\tiny] {n..n (la nivel de aplicație)};
\draw[->] (user) -- (role) node[midway, right, font=\tiny] {n..n};
\end{tikzpicture}
\caption{Diagrama entitate-relație simplificată a bazei de date CourseHub}
\label{fig:er-diagram}
\end{figure}

Un utilizator (\texttt{ApplicationUser}) poate fi atât expeditor, cât și destinatar al
mai multor mesaje, atât solicitant (elev), cât și țintă (profesor) a mai multor cereri
de colaborare, și poate avea mai multe coduri OTP și roluri asociate. Dintre acestea,
doar legătura cu \texttt{HireRequest} este materializată printr-o cheie externă
propriu-zisă; relația cu \texttt{OtpCode} este menținută la nivel de aplicație, prin
câmpul \texttt{UserId}.

Relația dintre \texttt{ApplicationUser} și \texttt{Topic} este, de asemenea, implementată
intenționat la nivel de aplicație, prin lista de identificatori din \texttt{SelectedTopics},
și nu printr-o tabelă de legătură: un utilizator poate selecta cel mult trei materii, iar
operațiile necesare sunt suficient de simple pentru a nu justifica o entitate suplimentară.

\section{Arhitectura logică a sistemului}

CourseHub este structurată modular, fiecare controller acoperind o arie de
responsabilitate clară:

\begin{itemize}
    \item \textbf{Identity (Areas/Identity)}: înregistrare, autentificare, verificare OTP,
    resetarea parolei, gestionarea profilului propriu (\texttt{Manage});
    \item \textbf{SearchController}: căutarea și filtrarea utilizatorilor;
    \item \textbf{HireController}: trimiterea, acceptarea, anularea cererilor de colaborare;
    \item \textbf{ConnectionsController}: afișarea listei de colaborări active (elevi, respectiv
    profesori, în funcție de rol);
    \item \textbf{ChatController} și \textbf{ChatHub}: afișarea istoricului unei conversații
    și, respectiv, livrarea mesajelor în timp real prin SignalR;
    \item \textbf{MessagesController}: lista mesajelor primite, cu conversațiile recente ale utilizatorului;
    \item \textbf{AdminController}: panoul de administrare (utilizatori, roluri, blocări,
    materii, moderarea mesajelor).
\end{itemize}

Figura~\ref{fig:arhitectura} ilustrează fluxul general al unei cereri prin sistem.

\begin{figure}[h]
\centering
\begin{tikzpicture}[
    box/.style={draw, rectangle, minimum width=3.4cm, minimum height=0.9cm, align=center, font=\small},
    >=Stealth
]
\node[box] (browser) {Interfață web\\(Razor Views, JS)};
\node[box, below=1.2cm of browser] (mvc) {ASP.NET Core\\(Controllere, Servicii)};
\node[box, below left=1.2cm and 2.2cm of mvc] (db) {SQL Server\\(Entity Framework Core)};
\node[box, below right=1.2cm and 2.2cm of mvc] (hub) {SignalR Hub\\(ChatHub)};

\draw[->] (browser) -- (mvc) node[midway, right, font=\tiny] {HTTP(S)};
\draw[->] (mvc) -- (db) node[midway, left, font=\tiny] {EF Core};
\draw[<->] (browser) -- (hub) node[midway, right, font=\tiny] {WebSocket};
\draw[->] (hub) -- (mvc) node[midway, below, font=\tiny] {};
\end{tikzpicture}
\caption{Arhitectura logică a platformei CourseHub}
\label{fig:arhitectura}
\end{figure}

\subsection{Fluxuri funcționale principale}

Parcursul unui utilizator leagă modulele de mai sus într-un singur lanț: se
înregistrează printr-un formular protejat de reCAPTCHA și își confirmă contul cu
un cod OTP, parcurge apoi fluxul de trei pași prin care își definește rolul,
detaliile și materiile de interes, caută un partener potrivit și îi trimite o
cerere de colaborare, iar din momentul acceptării cei doi comunică prin chat în
timp real, profesorul putând atașa suplimentar fișiere. În paralel,
administratorul gestionează conturile, rolurile, blocările și materiile
disponibile pe platformă.
